%----------------------------------------------------------------------------------------
\documentclass[a0,portrait]{a0poster}

\usepackage{poster_classe_nice_booklet}

\begin{document}

\onehalfspacing

%----------------------------------------------------------------------------------------
%	POSTER HEADER
%----------------------------------------------------------------------------------------

\vspace*{-7cm}
\begin{figure}[!htb]
   \begin{minipage}{0.65\textwidth}
     
\includegraphics[width=.85\linewidth]{LogoCNMAC2025}
   \end{minipage}\hfill
   \begin{minipage}{0.40\textwidth}
	 \hspace{12cm}
     
\includegraphics[width=.4\linewidth]{Logo_CEFET-MG}
   \end{minipage}
\end{figure}

\vspace{3mm}
\begin{minipage}[b]{1.0\linewidth} %default 0.75
\begin{center}
\Huge \textbf{Classe \LaTeX{} para Criação de e-Books para o PROFMAT}\\
\huge \textbf{Luis Alberto D'Afonseca}    \\[2mm]
\huge CEFET-MG, MG, Brasil                \\[2mm]
\Large \texttt{luis.dafonseca@cefetmg.br}
\end{center}
\end{minipage}
%
\vspace{-7mm}

%----------------------------------------------------------------------------------------

\begin{multicols}{2}
\raggedbottom

\justify

\begin{center}
  \centering
  \fbox{\includegraphics[width=0.875\linewidth]{NiceBooklet_Manual-capa}}%
  \hspace{-103mm}%
  \adjustbox{raise=405mm}{\fcolorbox{white}{white}{%
      \qrcode[height=100mm]{https://material-didatico.github.io/pages/booklet}}}
  \captionof{figure}{{
      Capa do manual, exemplo de uso, da classe \texttt{NiceBooklet}. \setlength\bibitemsep{0pt}
      Fonte: o próprio autor.
  }}
  \label{fig:capa}
\end{center}

%----------------------------------------------------------------------------------------
\vspace{-10mm}
\section*{Resumo}
\vspace{-10mm}

Apresentação de uma classe \LaTeX{} destinada a facilitar o desenvolvimento de
\textit{e-books} por alunos do PROFMAT.
O \LaTeX{} é a principal ferramenta para escrever textos matemáticos,
mas apresenta grandes desafios para iniciantes.
A classe criada esconde uma parte significativa da complexidade,
incluindo os principais pacotes matemáticos, para escrever em português e
inserir imagens, além de definir a formatação.
Além disso, a classe automatiza diversas tarefas, por exemplo, transferir as
respostas dos exercícios para um capítulo específico no final do texto.

\vspace{5mm}
\noindent
\textbf{Palavras-Chave:}
acessibilidade digital;
ensino;
material didático;
produção acadêmica.

%----------------------------------------------------------------------------------------
\vspace{-10mm}
\section*{A Classe \texttt{NiceBooket}}
\vspace{-5mm}

\noindent
\LaTeX{} é a ferramenta mais utilizada para escrever
textos matemáticos, devido, entre outros fatores, à sua qualidade tipográfica e
separação entre conteúdo e formatação.
Porém, ela é bastante intimidadora para principiantes.

\hspace{-5.5mm}%
\begin{minipage}[b]{0.55\linewidth}

  \justify

  \noindent
  Essa dificuldade se tornou especialmente evidente na orientação de
  dissertações para o PROFMAT, onde os mestrandos precisam atuar como
  professores de matemática durante a realização do curso.
  Desejando facilitar a criação de apostilas e que elas fossem
  visualmente atraentes para os alunos do Ensino Básico,
  desenvolvi a classe \LaTeX{} \texttt{NiceBooklet}.
  \vspace{13mm}

  \noindent
  Características gráficas:
  \begin{itemize}
    \setlength\itemsep{0em}
    \item Destaque de resultados importantes,
          Figs \ref{fig:pagina_teorema} e \ref{fig:pagina_exemplo}

    \item Fontes grandes (14pt) para facilitar a leitura em formato digital

    \item Formatação da capa a partir de uma imagem;

    \item Criação de \textit{links} de navegação em cada página;

    \item Deslocamento das respostas dos exercícios;

    \item Criação de \textit{links} com QR code;

    \item Criação de um índice remissivo.
  \end{itemize}

  \vspace{-10mm}
  \section*{Resultados}

  {
  Estas são algumas dissertações que criaram apostilas utilizando a classe:
  Gomes criou um material abordando a Matemática dos
  investimentos financeiros
  \cite{ANDERSON_LUIS_DE_AVILA_GOMES:CEFET:2024};
  %
  Vieira escreveu uma apostila para apresentar o conceito de
  infinito para alunos do Ensino Médio
  \cite{PIERRE_FERMAT_VIEIRA:CEFET:2023};
  %
  de Jesus desenvolveu uma apostila sobre a aplicação da
  Matemática no cálculo de materiais na Construção Civil
  \cite{VILMAR_PEREIRA_DE_JESUS:CEFET:2022};
  %
  \unskip\parfillskip 0pt \par
  }
\end{minipage}%
\hfill%
\begin{minipage}[b]{0.39\linewidth}
\begin{center}
  \fbox{\includegraphics[width=\linewidth]{Pagina_exemplo_1}}
  \captionof{figure}{
    Um teorema e sua demonstração.
    Fonte: Vieira \cite{PIERRE_FERMAT_VIEIRA:CEFET:2023}.
  }
  \label{fig:pagina_teorema}
  \vspace{2mm}

  \fbox{\includegraphics[width=\linewidth]{Pagina_exemplo_2}}
  \captionof{figure}{
    Uma definição e um exemplo.
    Fonte: Reis \cite{TAMISSON_SANTOS_REIS:CEFET:2022}.
  }
  \label{fig:pagina_exemplo}
\end{center}
\end{minipage}

%----------------------------------------------------------------------------------------

\vspace{3mm}
\noindent
%
Reis criou um jogo de RPG (Role-Playing Game) para ser
usado como ferramenta para o ensino de probabilidades
\cite{TAMISSON_SANTOS_REIS:CEFET:2022};
%
Santos escreveu uma apostila para o ensino de noções de
programação e métodos numéricos para o Ensino Médio
\cite{Edson:CEFET:2024}.

%-------------------------------------------------------------------------------
\vspace{-15mm}
\section*{Agradecimentos}
\vspace{-10mm}

\noindent
O autor agradece ao CEFET-MG e à FAPEMIG
pelo apoio financeiro (Chamada FAPEMIG 16/2024, Processo PCE-00114-25).

%-------------------------------------------------------------------------------
\vspace{-10mm}
\printbibliography

\end{multicols}
%-------------------------------------------------------------------------------

\vfill
\begin{center}%
  \fbox{\includegraphics[width=0.19\textwidth]{capa_civil}}\hspace{-53mm}\adjustbox{raise=179mm}{\fcolorbox{white}{white}{\qrcode[height=50mm]{https://material-didatico.github.io/pages/civil}}}\hspace{45mm}%
  \fbox{\includegraphics[width=0.19\textwidth]{capa_infinito}}\hspace{-53mm}\adjustbox{raise=179mm}{\fcolorbox{white}{white}{\qrcode[height=50mm]{https://material-didatico.github.io/pages/infinito}}}\hspace{45mm}%
  \fbox{\includegraphics[width=0.19\textwidth]{capa_met_num}}\hspace{-53mm}\adjustbox{raise=179mm}{\fcolorbox{white}{white}{\qrcode[height=50mm]{https://material-didatico.github.io/pages/met_num_medio}}}\hspace{45mm}%
  \fbox{\includegraphics[width=0.19\textwidth]{capa_zipmath}}\hspace{-53mm}\adjustbox{raise=179mm}{\fcolorbox{white}{white}{\qrcode[height=50mm]{https://material-didatico.github.io/pages/zip_math}}}
\end{center}

\end{document}
%-------------------------------------------------------------------------------
